% ============================================================================
%  CHAPTER 2 — Cơ sở lý thuyết
% ============================================================================
\section{Cơ sở lý thuyết}
\label{sec:theory}

Chương này trình bày các cơ sở lý thuyết nền tảng cho đồ án, bao gồm bài toán trích xuất thông tin tài liệu (DIE), tổng quan về các công trình nghiên cứu liên quan, cơ chế hoạt động của nhận dạng ký tự quang học (\gls{ocr}), và chi tiết hai kiến trúc mô hình học sâu đại diện: \model{Donut} (OCR-free) và \model{LayoutXLM} (OCR-based). Cuối cùng, chương định nghĩa các độ đo đánh giá cùng quy ước xử lý đặc biệt đối với giá trị rỗng.

% ============================================================================
\subsection{Bài toán trích xuất thông tin tài liệu}
\label{subsec:kie_problem}

Trích xuất thông tin tài liệu (Document Information Extraction - DIE) hay trích xuất thông tin then chốt (Key Information Extraction - KIE) là bài toán nhận dạng và trích xuất các thực thể thông tin có cấu trúc dưới dạng cặp khóa-giá trị (key-value pairs) từ hình ảnh tài liệu bán cấu trúc (như biên lai, hóa đơn, tờ khai).

Đầu vào là một ảnh biên lai $I \in \mathbb{R}^{H \times W \times C}$. Đầu ra mong muốn là một tập hợp các thực thể thông tin: $Y = \{(K_1, V_1), (K_2, V_2), \dots, (K_M, V_M)\}$, trong đó mỗi $K_j$ tương ứng với một trường thông tin đích (ví dụ: \field{store\_name}, \field{date}, \field{total}, \field{address}) và $V_j$ là văn bản đã chuẩn hóa được trích xuất từ tài liệu. Quá trình này yêu cầu mô hình đồng thời nhận diện được ký tự văn bản và hiểu được mối quan hệ không gian, ngữ nghĩa giữa các vùng chữ để liên kết chúng một cách chính xác.

% ============================================================================
\subsection{Các công trình nghiên cứu liên quan}
\label{subsec:related_work}

Bài toán trích xuất thông tin tài liệu đã đạt được nhiều bước tiến lớn nhờ sự phát triển của các kiến trúc Transformer và mô hình đa phương thức. \cref{tab:method-comparison} so sánh tổng quan ba hướng tiếp cận được khảo sát trong đồ án cùng giả thuyết thực nghiệm đi kèm.

\begin{table}[htbp]
  \centering
  \caption{So sánh đặc tính và giả thuyết kiểm chứng của ba hướng tiếp cận.}
  \label{tab:method-comparison}
  \small
  \begin{tabularx}{\textwidth}{l p{2.8cm} L{3.5cm} L{4.5cm}}
    \toprule
    \textbf{Phương pháp} & \textbf{Đặc trưng kỹ thuật} & \textbf{Ưu điểm \& Hạn chế} & \textbf{Giả thuyết kiểm chứng} \\
    \midrule
    \textbf{Baseline (Regex)} 
      & Đối sánh từ khóa và biểu thức chính quy (Regex) trên kết quả OCR. 
      & Tốc độ xử lý cực nhanh, không cần dữ liệu huấn luyện. Tuy nhiên, luật heuristics cứng nhắc, dễ thất bại trước bố cục đa dạng.
      & Hoạt động hiệu quả trên các trường có định dạng cấu trúc cố định (\field{date}, \field{total}); kém hiệu quả trên các trường tự do (\field{store\_name}, \field{address}). \\
    \addlinespace
    \textbf{\model{LayoutXLM} (OCR-based)}
      & Phân loại token (\gls{bio} tagging) đa phương thức: Text + Layout (2D BBox) + Visual.
      & Độ chính xác cao, khai thác sâu thông tin không gian. Tuy nhiên, pipeline phức tạp, phụ thuộc trực tiếp vào chất lượng OCR đầu vào.
      & Đạt hiệu năng tối ưu nhất nhờ cơ chế nhúng không gian 2D, nhưng dễ bị đứt gãy gộp thực thể đa dòng khi OCR bị lỗi sắp xếp thứ tự đọc. \\
    \addlinespace
    \textbf{\model{Donut} (OCR-free)}
      & Mô hình sinh (generative) end-to-end từ ảnh sang văn bản có cấu trúc (thẻ XML).
      & Loại bỏ hoàn toàn bộ OCR độc lập, giảm thiểu tích lũy sai số. Tuy nhiên, tốc độ suy luận chậm, dễ bị ảo giác (hallucination).
      & Khó hội tụ và dễ mắc lỗi sinh ảo hoặc lỗi định dạng thẻ tag khi kích thước tập dữ liệu huấn luyện nhỏ (khoảng 1.000 mẫu). \\
    \bottomrule
  \end{tabularx}
\end{table}

Các nghiên cứu nền tảng và công trình liên quan trực tiếp đến đồ án bao gồm:
\begin{itemize}
  \item \textbf{MC-OCR 2021:} Cuộc thi xử lý biên lai chụp di động do Zalo AI tổ chức~\cite{mcocr2021}, cung cấp bộ dữ liệu biên lai tiếng Việt thực tế làm cơ sở cho đồ án. Các giải pháp hàng đầu trong cuộc thi thường kết hợp bộ phát hiện chữ (DBNet, YOLO) và bộ nhận dạng mạnh mẽ, kết hợp với các luật hậu xử lý sâu.
  \item \textbf{CORD:} Bộ dữ liệu biên lai quốc tế Consolidated Receipt Dataset~\cite{park2019cord}, đóng vai trò benchmark chính cho các kiến trúc hiểu tài liệu end-to-end.
  \item \textbf{LayoutLM và LayoutLMv3:} Các mô hình tiền huấn luyện đa phương thức của Microsoft~\cite{xu2020layoutlm, huang2022layoutlmv3}. \model{LayoutXLM}~\cite{xu2021layoutxlm} là phiên bản mở rộng đa ngôn ngữ dựa trên XLM-RoBERTa, đóng vai trò cốt lõi trong pipeline OCR-based tiếng Việt của đồ án.
  \item \textbf{TrOCR:} Mô hình nhận dạng chữ viết (Text Recognition) dựa hoàn toàn trên kiến trúc Transformer Encoder-Decoder~\cite{li2023trocr}. TrOCR nhận đầu vào là một ảnh dòng chữ được cắt ra từ bộ phát hiện chữ viết độc lập (Text Detector) và giải mã thành văn bản, đạt hiệu năng rất cao nhờ cơ chế tự chú ý.
\end{itemize}

% ============================================================================
\subsection{Nhận dạng ký tự quang học (OCR)}
\label{subsec:ocr}

Hệ thống OCR hiện đại thường được thiết kế thành một pipeline gồm hai giai đoạn độc lập: phát hiện văn bản và nhận dạng chữ viết. Trong pipeline OCR-based của đồ án, module OCR đảm nhận vai trò trích xuất văn bản thô cùng tọa độ 2D từ ảnh biên lai để làm đầu vào cho mô hình \model{LayoutXLM}.

\subsubsection{Phát hiện văn bản với PaddleOCR}
\label{subsubsec:paddleocr}

\model{PaddleOCR}~\cite{du2020ppocr} sử dụng thuật toán Differentiable Binarization (DB) để phát hiện các vùng chứa chữ viết. Khác với các thuật toán nhị phân hóa truyền thống sử dụng ngưỡng cố định, DBNet tích hợp một bước nhị phân hóa khả vi (differentiable binarization) vào mạng nơ-ron:
\begin{equation}
  \hat{B}_{i,j} = \frac{1}{1 + \exp(-k(P_{i,j} - T_{i,j}))}
  \label{eq:db_binarization}
\end{equation}
trong đó $P_{i,j}$ là bản bản đồ xác suất chứa ký tự, $T_{i,j}$ là bản đồ ngưỡng thích ứng học được từ mạng backbone và FPN, và $k$ là hệ số phóng đại khả vi. Cơ chế này cho phép tối ưu hóa trực tiếp ranh giới hộp bao văn bản thông qua lan truyền ngược, đem lại độ chính xác cao đối với các dòng chữ nhỏ và nghiêng trên biên lai.

\subsubsection{Nhận dạng văn bản tiếng Việt với VietOCR}
\label{subsubsec:vietocr}

Ảnh của các dòng chữ sau khi được phát hiện sẽ được đưa vào \model{VietOCR}~\cite{vietocr2020} để nhận dạng chuỗi ký tự tiếng Việt. VietOCR kết hợp một mạng CNN (VGG-19 được tinh chỉnh) để trích xuất đặc trưng thị giác từ ảnh dòng chữ, và một bộ giải mã Transformer Encoder-Decoder tuần tự để giải mã các đặc trưng này thành chuỗi Unicode tiếng Việt có dấu. 

Mô hình được huấn luyện bằng hàm mất mát Connectionist Temporal Classification (CTC)~\cite{graves2006ctc}, cho phép huấn luyện trực tiếp trên các cặp ảnh-văn bản mà không cần căn chỉnh vị trí từng ký tự ở mức pixel:
\begin{equation}
  \mathcal{L}_{\text{CTC}} = -\ln P(\mathbf{y} \mid \mathbf{x}) = -\ln \sum_{\mathbf{\pi} \in \mathcal{B}^{-1}(\mathbf{y})} P(\mathbf{\pi} \mid \mathbf{x})
  \label{eq:ctc_loss}
\end{equation}
trong đó $\mathbf{x}$ là chuỗi đặc trưng visual, $\mathbf{y}$ là chuỗi ký tự đích, và $\mathcal{B}$ là hàm rút gọn loại bỏ ký tự trống (blank) và các ký tự lặp liên tiếp.

% ============================================================================
\subsection{Mô hình Donut (OCR-free)}
\label{subsec:donut}

\model{Donut} (Document Understanding Transformer)~\cite{kim2022donut} là kiến trúc Transformer end-to-end ánh xạ trực tiếp điểm ảnh của tài liệu thành văn bản có cấu trúc dưới dạng thẻ tag XML mà không cần thông qua module OCR trung gian.

\subsubsection{Kiến trúc tổng quan}
\label{subsubsec:donut_arch}

Kiến trúc Donut bao gồm hai thành phần chính:
\begin{enumerate}
  \item \textbf{Vision Encoder (Swin Transformer):} Mã hóa hình ảnh biên lai thành chuỗi biểu diễn đặc trưng thị giác phân cấp $\mathbf{z} = \{z_1, z_2, \ldots, z_n\}$. Swin Transformer~\cite{liu2021swin} sử dụng cơ chế cửa sổ tự chú ý trượt (Shifted Window MSA) để giảm độ phức tạp tính toán từ bậc hai xuống bậc tuyến tính theo kích thước ảnh, đồng thời duy trì sự liên tục ngữ cảnh.
  \item \textbf{Text Decoder (mBART):} Bộ giải mã tự hồi quy (autoregressive) đa ngôn ngữ~\cite{lewis2020bart}. Tại mỗi bước thời gian $i$, decoder kết hợp thông tin tự chú ý của các token đã sinh phía trước và thực hiện cơ chế Cross-Attention lên biểu diễn đặc trưng thị giác $\mathbf{z}$ của Encoder để dự đoán xác suất token tiếp theo:
  \begin{equation}
    \mathcal{L}_{\text{Donut}} = -\sum_{i=1}^{T} \log P(y_i \mid y_{<i}^*, \mathbf{z})
    \label{eq:donut_loss}
  \end{equation}
\end{enumerate}

\begin{figure}[htbp]
  \centering
  \includegraphics[width=0.85\textwidth]{donut_architecture.jpg}
  \caption{Kiến trúc tổng quan mô hình \model{Donut} (Swin Encoder kết hợp mBART Decoder)~\cite{kim2022donut}.}
  \label{fig:donut_arch}
\end{figure}

\subsubsection{Phân biệt khái niệm Warm-up trong Donut}
\label{subsubsec:donut_warmup}

Trong quá trình huấn luyện mô hình Donut, cần phân biệt rõ hai khái niệm "warm-up" dễ gây nhầm lẫn:
\begin{itemize}
  \item \textbf{Learning Rate Warm-up (tỷ lệ $warmup\_ratio$ hoặc số bước $warmup\_steps$):} Đây là một kỹ thuật điều tiết tốc độ học trong tối ưu hóa (LR scheduling). Trong số bước đầu tiên (thường là 10\% tổng số bước huấn luyện), tốc độ học được tăng dần tuyến tính từ 0 lên mức cực đại được thiết lập nhằm tránh hiện tượng bùng nổ gradient khi bắt đầu cập nhật trọng số.
  \item \textbf{Pre-training/Warm-up mô hình trên tập CORD:} Đây là một bước huấn luyện phụ (pre-finetuning). Trước khi huấn luyện trên tập dữ liệu biên lai tiếng Việt chính, mô hình Donut được huấn luyện trên một tập dữ liệu biên lai lớn khác (như CORD v2 gồm 800 mẫu tiếng Anh) để mô hình làm quen với bố cục và cấu trúc biểu diễn thẻ XML tổng quát của biên lai.
\end{itemize}

% ============================================================================
\subsection{Mô hình LayoutXLM (OCR-based)}
\label{subsec:layoutxlm}

\model{LayoutXLM}~\cite{xu2021layoutxlm} là mô hình đa phương thức đa ngôn ngữ tiền huấn luyện dành riêng cho các tài liệu có cấu trúc bố cục phong phú, kế thừa backbone XLM-RoBERTa~\cite{conneau2020xlmr}.

\subsubsection{Nhúng đa phương thức (Multimodal Embedding)}
\label{subsubsec:multimodal_embedding}

Để biểu diễn một token văn bản trong tài liệu, LayoutXLM kết hợp ba loại vector nhúng:
\begin{enumerate}
  \item \textbf{Nhúng văn bản (Text Embedding):} Mã hóa đặc trưng ngữ nghĩa của từ.
  \item \textbf{Nhúng không gian 2D (Spatial Embedding):} Tọa độ bounding box $[x_1, y_1, x_2, y_2]$ của từ được chuẩn hóa về khoảng $[0, 1000]$ và nhúng thông qua các bảng nhúng tọa độ để mô hình nắm bắt vị trí tương đối và tuyệt đối trên ảnh:
  \begin{equation}
    \mathbf{e}_{2D}(\text{bbox}_i) = \mathbf{e}_{x_1}(x_1) + \mathbf{e}_{y_1}(y_1) + \mathbf{e}_{x_2}(x_2) + \mathbf{e}_{y_2}(y_2) + \mathbf{e}_{w}(w) + \mathbf{e}_{h}(h)
  \end{equation}
  \item \textbf{Nhúng thị giác (Visual Embedding):} Cắt vùng ảnh tương ứng với bounding box của token và đưa qua mạng CNN (ResNeXt) để trích xuất đặc trưng hình ảnh (phông chữ, độ đậm nhạt, màu sắc).
\end{enumerate}

Vector biểu diễn tổng hợp của token là tổng của ba thành phần trên: $\mathbf{x}_i = \mathbf{t}_i + \mathbf{e}_{2D}(\text{bbox}_i) + \mathbf{v}_i$.

\subsubsection{Spatial-aware Self-Attention}
\label{subsubsec:spatial_attention}

Cơ chế tự chú ý nhận biết không gian (spatial-aware) của LayoutXLM bổ sung thêm các thông tin thiên lệch vị trí (bias) theo chiều dọc và chiều ngang dựa trên khoảng cách vật lý giữa hai bounding box của token $i$ và token $j$:
\begin{equation}
  \alpha_{ij} = \frac{\mathbf{q}_i \mathbf{k}_j^T}{\sqrt{d}} + b^{1D}(j - i) + b^{2D}_x(x_{1,i} - x_{1,j}) + b^{2D}_y(y_{1,i} - y_{1,j})
\end{equation}
Điều này giúp mô hình tự động học được mối liên hệ logic giữa các từ nằm gần nhau về mặt không gian địa lý trên ảnh tài liệu.

\subsubsection{Gán nhãn BIO (BIO Tagging)}
\label{subsubsec:bio_tagging}

Để thực hiện trích xuất thông tin, bài toán được định nghĩa dưới dạng phân loại token (token classification). Mỗi token đầu vào được dự đoán một nhãn trong hệ thống nhãn \gls{bio} gồm 9 nhãn được mô tả chi tiết trong \cref{tab:bio_labels}.

\begin{table}[htbp]
  \centering
  \caption{Hệ thống nhãn \gls{bio} cho 4 trường thông tin trích xuất.}
  \label{tab:bio_labels}
  \begin{tabular}{lll}
    \toprule
    \textbf{Trường thông tin} & \textbf{Nhãn Bắt đầu (Begin)} & \textbf{Nhãn Bên trong (Inside)} \\
    \midrule
    Tên cửa hàng (\field{store\_name}) & \texttt{B-STORE\_NAME} & \texttt{I-STORE\_NAME} \\
    Ngày tháng (\field{date})           & \texttt{B-DATE}        & \texttt{I-DATE}        \\
    Tổng tiền (\field{total})           & \texttt{B-TOTAL}       & \texttt{I-TOTAL}       \\
    Địa chỉ (\field{address})           & \texttt{B-ADDRESS}     & \texttt{I-ADDRESS}     \\
    Trường không quan tâm               & \multicolumn{2}{l}{\texttt{O} (Outside)} \\
    \bottomrule
  \end{tabular}
\end{table}

% ============================================================================
\subsection{Các độ đo đánh giá và quy ước xử lý rỗng}
\label{subsec:metrics}

Để đánh giá chính xác chất lượng trích xuất, các dự đoán sau khi được chuẩn hóa sẽ được so sánh với nhãn ground-truth thông qua ba chỉ số chính: EM, NES, và CER. 

\subsubsection{Exact Match (EM)}
EM đo tỷ lệ khớp chính xác hoàn toàn từng ký tự giữa dự đoán ($pred$) và nhãn thực ($gold$):
\begin{equation}
  \metric{EM}(pred, gold) = \begin{cases} 
    1.0 & \text{nếu } pred.strip() = gold.strip() \\
    0.0 & \text{ngược lại}
  \end{cases}
\end{equation}
\textbf{Quy ước giá trị rỗng:} Nếu cả $pred$ và $gold$ đều là chuỗi rỗng sau chuẩn hóa $\to \metric{EM} = 1.0$. Nếu chỉ một trong hai chuỗi rỗng $\to \metric{EM} = 0.0$.

\subsubsection{Normalized Edit Similarity (NES)}
NES đo mức độ tương đồng ngữ nghĩa ký tự thông qua khoảng cách Levenshtein chuẩn hóa:
\begin{equation}
  \metric{NES}(pred, gold) = 1 - \frac{\text{Levenshtein}(pred, gold)}{\max(\text{len}(pred), \text{len}(gold))}
\end{equation}
\textbf{Quy ước giá trị rỗng:} Nếu cả $pred$ và $gold$ đều là chuỗi rỗng $\to \metric{NES} = 1.0$ (tránh lỗi chia cho 0). Nếu chỉ một trong hai chuỗi rỗng $\to \metric{NES} = 0.0$.

\subsubsection{Character Error Rate (CER)}
CER đo tỷ lệ lỗi ở cấp độ ký tự trên chuỗi nhãn thực:
\begin{equation}
  \metric{CER}(pred, gold) = \frac{\text{Levenshtein}(pred, gold)}{\text{len}(gold)}
\end{equation}
\textbf{Quy ước giá trị rỗng:} 
\begin{itemize}
  \item Nếu cả hai cùng rỗng $\to \metric{CER} = 0.0$.
  \item Nếu $gold$ rỗng nhưng $pred$ có giá trị $\to \metric{CER} = 1.0$.
  \item Nếu $gold$ có giá trị nhưng $pred$ rỗng $\to \metric{CER} = 1.0$.
\end{itemize}
\textbf{Giải thích hiện tượng CER > 100\%:} Chỉ số CER có thể vượt quá 100\% (hay > 1.0) khi chuỗi dự đoán dài hơn đáng kể so với nhãn thực do mô hình sinh thừa nhiều ký tự rác hoặc bị lỗi lặp chuỗi tự hồi quy (hallucination), dẫn đến khoảng cách Levenshtein lớn hơn độ dài ký tự của nhãn thực $gold$.
